\section{Conclusion}\label{sec:conclusion}

A hierarchical latent-class model of reference agreement can be specified, fitted to convergence and checked without a gold standard, and its claim tiers are enriched for held-out evidence.
It does not, however, give a better ranking than simple combination rules; on the ChIP references the main reason is that the assumption it shares with every standard estimator, that a pair no resource covers is missing at random, is false.
Its error statements, moreover, are only as good as the selected posterior mode, because literature consensus and literature copying are only weakly distinguishable without independent evidence modalities.
For GRN benchmarking the practical value of the framework lies less in the estimator than in the evaluation design around it: separating the evidence that defines credibility from the evidence that tests it, reporting prevalence-free and TF-stratified metrics, controlling the posterior expected error rate of claims, and stating the conditions under which those claims hold.
